% Part: propositional-logic
% Chapter: syntax-and-semantics
% Section: semantic-notions

\documentclass[../../../include/open-logic-section]{subfiles}

\begin{document}

\olfileid{pl}{syn}{sem}

\olsection{Semantische begrippen}

We definiëren nu de volgende begrippen over waarheid en betekenis:

\begin{defn} 
\begin{enumerate}
\item !!^a{formula}~$!A$ is \emph{vervulbaar}, dus waar te maken, als
  er een~$\pAssign{v}$ bestaat waarvoor $\pSat{v}{!A}$. De formule is
  \emph{onvervulbaar} als er geen $\pAssign{v}$ bestaat waarvoor
  $\pSat{v}{!A}$;
\item !!^a{formula}~$!A$ is een \emph{tautologie} als
  $\pSat{v}{!A}$ voor alle !!{valuation}s~$\pAssign{v}$;
\item !!^a{formula}~$!A$ is \emph{contingent} als zij wel vervulbaar
  is, maar geen tautologie;
\item Laat $\Gamma$ een verzameling !!{formula}s zijn. Dan geldt
  $\Gamma \Entails !A$ (``uit $\Gamma$ volgt $!A$'') precies als
  $\pSat{v}{!A}$ geldt voor elke !!{valuation}~$\pAssign{v}$ waarvoor
  $\pSat{v}{\Gamma}$;
\item Laat $\Gamma$ een verzameling !!{formula}s zijn. De verzameling
  $\Gamma$ is \emph{vervulbaar} als er !!a{valuation}~$\pAssign{v}$
  bestaat waarvoor $\pSat{v}{\Gamma}$. Anders is $\Gamma$
  \emph{onvervulbaar}.
\end{enumerate} 
\end{defn}

\begin{prob} 
 Bekijk de volgende vier !!{formula}s. Bepaal voor elke formule of zij
 (a)~vervulbaar, (b)~een tautologie en (c)~contingent is.
\begin{enumerate}
  \item \( ( \Obj p_0 \lif ( \lnot \Obj p_1 \lif \lnot \Obj p_0 ) ) \).
  \item \( ( ( \Obj p_0 \land \lnot \Obj p_1 ) \lif ( \lnot \Obj p_0 \land \Obj p_2 )) \liff ( ( \Obj p_2 \lif \Obj p_0 ) \lif ( \Obj p_0 \lif \Obj p_1 )) \).
  \item \( ( \Obj p_0 \liff \Obj p_1 ) \lif ( \Obj p_2 \liff \lnot \Obj p_1 ) \).
  \item \( (( \Obj p_0 \liff ( \lnot \Obj p_1 \land \Obj p_2 )) \lor ( \Obj p_2 \lif ( \Obj p_0 \liff \Obj p_1 ))) \).
\end{enumerate}
\end{prob}

\begin{prop}
\ollabel{prop:semanticalfacts} 
\begin{enumerate} 
\item $!A$ is een tautologie precies als
  $\emptyset \Entails !A$; 
\item Als $\Gamma \Entails !A$ en $\Gamma \Entails !A \lif !B$, dan
  geldt $\Gamma \Entails !B$;
\item Als $\Gamma$ vervulbaar is, dan is elke eindige deelverzameling
  van $\Gamma$ ook vervulbaar; 
\item \ollabel{def:monotonicity} Monotoniciteit: als
  $\Gamma \subseteq \Delta$ en $\Gamma \Entails !A$, dan geldt ook
  $\Delta \Entails !A$;
\item \ollabel{def:Cut} Transitiviteit: als $\Gamma \Entails !A$ en
  $\Delta \cup \{ !A\} \Entails !B$, dan geldt
  $\Gamma \cup \Delta \Entails !B$.
\end{enumerate}
\end{prop}

\begin{proof}
Opgave.
\end{proof}

\begin{prob}
Bewijs \olref[pl][syn][sem]{prop:semanticalfacts}
\end{prob}

\begin{prop}\ollabel{prop:entails-unsat}
  $\Gamma \Entails !A$ precies als $\Gamma \cup \{\lnot !A\}$
  onvervulbaar is.
\end{prop}

\begin{proof}
Opgave.
\end{proof}

\begin{prob}
Bewijs \olref[pl][syn][sem]{prop:entails-unsat}
\end{prob}

\begin{thm}[Semantische deductiestelling]
  \ollabel{thm:sem-deduction} $\Gamma \Entails !A \lif !B$ precies als
  $\Gamma \cup \{!A\} \Entails !B$.
\end{thm}

\begin{proof}
Opgave.
\end{proof}

\begin{prob}
Bewijs \olref[pl][syn][sem]{thm:sem-deduction}
\end{prob}
\end{document}
